\documentclass[12pt]{article}
\usepackage[spanish]{babel}
\usepackage[utf8]{inputenc}
\usepackage[T1]{fontenc}
\usepackage{geometry}
\usepackage{graphicx}
\usepackage{setspace}
\usepackage{titlesec}
\usepackage{hyperref}

\geometry{margin=1in}
\onehalfspacing

\begin{document}

\begin{titlepage}
    \centering

    % Logo opcional
    %\includegraphics[width=0.3\textwidth]{logo.png}\par\vspace{1cm}

    {\scshape\Large Inteligencia Artificial II \par}
    \vspace{1.5cm}

    {\huge\bfseries Memoria de trabajo \par}
    \vspace{0.5cm}
    {\Large Curso de Verano \par}

    \vfill

    {\large
    \begin{tabular}{rl}
        Nombre del Alumno: & Emilio Izquierdo Montero \\
        Profesor:          & Carlos Alberto Sanchez Garcia \\
        Grupo / Semestre:  & 9 \\
        Fecha:             & 30 de Junio del 2026 \\
    \end{tabular}
    \par}

    \vfill

\end{titlepage}

\tableofcontents
\newpage

\section{Introducción}

La presente memoria técnica documenta las actividades, prácticas y proyectos desarrollados durante el curso de verano de la asignatura 
Inteligencia Artificial II, impartida por el profesor Carlos.
En este documento se recopilan los conocimientos, herramientas y metodologías estudiadas a lo largo de la materia, 
haciendo énfasis en el desarrollo e implementación de modelos de aprendizaje automático mediante el uso de TensorFlow.
Asimismo, se presentan los resultados obtenidos y las experiencias adquiridas durante el proceso de aprendizaje.

\newpage

\section{Objetivo}

Aplicar los fundamentos de la inteligencia artificial mediante el diseño, entrenamiento y evaluación de modelos de aprendizaje automático 
utilizando TensorFlow, con el fin de comprender su funcionamiento y desarrollar soluciones capaces de resolver problemas de clasificación,
predicción y reconocimiento de patrones.

\newpage

\section{Unidad 1}
\subsection{¿Qué es TensorFlow?}

TensorFlow es una biblioteca de código abierto especializada en aprendizaje automático e inteligencia artificial, desarrollada originalmente por la empresa Google, 
Esta plataforma proporciona herramientas para diseñar, entrenar y desplegar modelos de redes neuronales y otros algoritmos de aprendizaje automático, 
permitiendo aprovechar tanto procesadores convencionales (CPU) como aceleradores especializados como GPU y TPU.

\vspace{0.5cm}

Gracias a su flexibilidad y escalabilidad, TensorFlow es ampliamente utilizado en entornos académicos, de investigación e industriales para resolver problemas de clasificación, 
regresión, visión por computadora, procesamiento de lenguaje natural y análisis de datos

\vspace{0.5cm}

El proyecto cuenta con una de las comunidades de desarrollo más grandes dentro del ecosistema de inteligencia artificial de código abierto. 
Actualmente participan más de 3,800 colaboradores provenientes de distintas organizaciones y países.
TensorFlow está desarrollado principalmente en C++ (56.1\%) y Python (25.3\%).


\vspace{0.5cm}
Para más información, consulta la \href{https://www.tensorflow.org}{página oficial de TensorFlow}.
\newpage
\subsection{Instalación de tensorflow}
Para trabajar con TensorFlow de manera local fue necesario crear un entorno virtual de Python, un entorno virtual permite aislar las dependencias de un proyecto del resto del sistema, evitando conflictos entre versiones de bibliotecas y facilitando la administración de paquetes, escencial en sistemas linux.


La creación del entorno virtual se realizó mediante el siguiente comando:

\begin{verbatim}
python -m venv venv
\end{verbatim}

Posteriormente, se activó el entorno virtual para que las instalaciones de paquetes se realizaran únicamente dentro de este entorno:

\begin{verbatim}
source venv/bin/activate
\end{verbatim}

Una vez activado el entorno, se procedió a instalar TensorFlow utilizando el gestor de paquetes \texttt{pip}:

\begin{verbatim}
pip install tensorflow
\end{verbatim}

\begin{figure}[h]
    \centering
    \includegraphics[width=0.9\linewidth]{/home/emilio/Imágenes/memoriaia/InstalacionTensorflow.png}
    \caption{Instalacion de tensorflow con pip}
\end{figure}

Una ves hecho esto verificamos la lista de paquetes instalados con

\begin{verbatim}
pip list
\end{verbatim}

\begin{figure}[h]
    \centering
    \includegraphics[width=0.4\linewidth]{/home/emilio/Imágenes/memoriaia/listpip.png}
    \caption{Paquetes de pip}
\end{figure}

\subsection{Actividad 2: Perceptrones NADN, AND OR}

Vemos que con modelos neuronales sencilos podemos crear perceptones que tomen decisicones como las computadoras usando compuertas logicas como XOR AND 


\subsection{Actividad 3: Perceptrone XOR}

\subsection{Actividad 4: Iris Dataser perceptron}

\subsection{Actividad 5: Gradiente Descendiente y perceptrones}
\subsection{Actividad 6: Gradiente Descendiente y una red neuronal artificial}

\subsection{Actividad 6: Gradiente Descendiente y una red neuronal artificial}
\subsection{Actividad 7: Forward Pass y BackPropagation en Red Neuronal Artificial}
\subsection{Actividad 8: Celebrities Dataset TensorFlow}
\subsection{Actividad 9: Regresion lineal y clasificación con keras}
\subsection{Actividad 10: Salaru prediction con TensorFlow}
\subsection{Actividad 11: Resumen de modelos de entrenamiento tabla 1}
\subsection{Actividad 12: entrenamiento de modelo optimizado y GUI}
\subsection{Actividad 13: Clasificación binaria student drop out}
\subsection{Actividad 14: Clasificación con dataset employees}
\subsection{Actividad 15: Handwriting Classification model}
\subsection{Actividad 16: Mejorando el modelo base del ejercicio MNIST fashion}
\subsection{Actividad 17: Weather Classification}
\subsection{Actividad 18: Carta tecnica de transfer learging tabla 2}
\subsection{Actividad 19: Carta tecnica 8 modelos handling overtfitting}
\subsection{Actividad 20: The yelp polarity datasets}
\subsection{Actividad 21: modelos para evitar el overtfitting yelp}
\subsection{Actividad 22: Naive Forecast con dataset mercantil}
\subsection{Actividad 23: Sales forescasting using neural network}
\subsection{Conclusion y pensamientos finales}
\end{document}
